% Gerado por tcc_dissertation/scripts/gerar_tabelas.py a partir de resultados/. Não edite à mão.
\begin{table}[!htbp]
\centering
\caption{\EP{Composition of the five test folds (\texttt{StratifiedGroupKFold}, seed 42, groups = case, stratified by the four-class label). Contact: images (cases). No case appears in more than one fold.}{Composição das cinco dobras de teste (\texttt{StratifiedGroupKFold}, semente 42, grupos = caso, estratificação pelo rótulo de quatro classes). Encostado: imagens (casos). Nenhum caso aparece em mais de uma dobra.}}
\label{tab:dobras}
\footnotesize
\setlength{\tabcolsep}{4pt}
\begin{tabular}{crrrrrr}
\toprule
\EP{Fold}{Dobra} & \EP{Images}{Imagens} & \EP{Cases}{Casos} & \EP{Exits}{Saídas} & \EP{Contact}{Encostado} & \EP{Close}{Curta} & \EP{Distant}{Distância} \\
\midrule
1 & 506 & 96 & 146 & 4 (2) & 53 & 303 \\
2 & 459 & 95 & 118 & 10 (4) & 25 & 306 \\
3 & 531 & 99 & 133 & 13 (5) & 43 & 342 \\
4 & 538 & 100 & 132 & 11 (3) & 61 & 334 \\
5 & 520 & 96 & 142 & 4 (2) & 50 & 324 \\
\midrule
\EP{Total}{Total} & 2{,}554 & 486 & 671 & 42 (16) & 232 & 1{,}609 \\
\bottomrule
\end{tabular}
\end{table}
